% Exercise 3.2-3

\begin{proof}
  Since $n! \leq n^n$ for all $n \geq 1$, we have $\lg{(n!)} \leq
  \lg{n^n} = n\lg{n}$
  for all $n \geq 1$, so that we have $\lg{(n!)} \leq c_2 \cdot (n\lg{n})$ for
  $n \geq n_0$ with the identifications $c_2 = 1$ and $n_0 = 1$, i.e.
  $\lg{(n!)} = O(n\lg{n})$.

  For the lower bound we cannot take $c_1 = 1$, since $\lg{(n!)} < n\lg{n}$ for
  every $n \geq 2$; we must find some constant $c_1 < 1$ instead. From
  Stirling's approximation,
  \[
    n! = \sqrt{2\pi n} {\biggl(\frac{n}{e}\biggr)}^n
    \biggl(1 + \Theta \biggl(\frac{1}{n}\biggr)\biggr),
  \]
  and taking logarithms,
  \begin{equation*}
    \begin{split}
      \lg{(n!)} &= \frac{1}{2}\lg{(2\pi n)} + n\lg{n} - n\lg{e}
      + \lg{\biggl(1 + \Theta\biggl(\frac{1}{n}\biggr)\biggr)} \\
      &= n\lg{n} - n\lg{e} + \frac{1}{2}\lg{n} + \frac{1}{2}\lg{(2\pi)}
      + \Theta\biggl(\frac{1}{n}\biggr).
    \end{split}
  \end{equation*}
  The terms $\frac{1}{2}\lg{n}$ and $\frac{1}{2}\lg{(2\pi)}$ are positive for
  $n \geq 1$, so discarding them can only weaken the bound, and it suffices to
  show that $n\lg{n} - n\lg{e} \geq \frac{1}{2}n\lg{n}$. This holds exactly when
  \begin{equation*}
    \begin{split}
      \frac{1}{2}n\lg{n} &\geq n\lg{e} \\
      \Leftrightarrow \lg{n} &\geq 2\lg{e} \approx 2.885 \\
      \Leftrightarrow n &\geq e^2 \approx 7.39,
    \end{split}
  \end{equation*}
  so that $\lg{(n!)} \geq \frac{1}{2}(n\lg{n})$ for all $n \geq 8$, and
  $\lg{(n!)} = \Omega(n\lg{n})$ with $c_1 = \frac{1}{2}$ and $n_0 = 8$.

  Since $\lg{(n!)}$ is both $O(n\lg{n})$ and $\Omega(n\lg{n})$, exercise 3.1-5
  gives $\lg{(n!)} = \Theta(n\lg{n})$.

  Now,
  \[
    \frac{n!}{n^n} = \prod_{k = 1}^n \frac{k}{n}
    = \frac{1}{n} \cdot \prod_{k = 2}^n \frac{k}{n} \leq \frac{1}{n},
  \]
  since every factor with $k \leq n$ is at most 1. Hence $n!/n^n \to 0$ as
  $n \to \infty$, and $n! = o(n^n)$.

  Finally, writing $r_n = n!/2^n$, we have $r_n = r_{n - 1} \cdot \frac{n}{2}$,
  and $\frac{n}{2} \geq 2$ for all $n \geq 4$. Thus $r_n \geq 2r_{n-1}$ from
  $n = 4$ onwards, so $r_n \to \infty$ and $n! = \omega(2^n)$.
\end{proof}
